% GENERATED FILE. Edit canonical lessons, then run npm run generate:books.
% canonical-source-hash: fnv1a64:e141733e804edc39
% canonical-lessons: MR-C333-mohim, MR-C333-svari, MR-C333-maghar-ghene, MR-C333-mukta-karne, MR-C333-hakalne

\chapter{Campaigns}
\label{ch:mr-campaigns}
\hlchaptermodality{333}

\begin{chapteropening}
\textbf{By the end of this chapter:} \emph{I can tell of a campaign, a retreat and a release.}

Complete the last lesson of chapter 333: I can tell of a campaign, a retreat and a release.
\end{chapteropening}

\section[\texorpdfstring{\mr{मोहीम} (mohīm)}{मोहीम (mohīm)}]{\mr{मोहीम} (mohīm) — an expedition; a campaign}

\label{lesson:MR-C333-mohim}



\begin{quote}
Before the new one: say the Marathi for a sacrifice, then the Marathi for to surrender.
\end{quote}

\subsection*{You'll want to know: \mr{मोहीम}}
\textbf{\mr{मोहीम}} — \emph{mohīm} — \textquotedblleft{}an expedition; a campaign\textquotedblright{}.

\begin{quote}
\mr{गिर्यारोहकांची} \mr{मोहीम} \mr{यशस्वी} \mr{झाली}.
\end{quote}

\begin{quote}
\emph{giryārohkāñcī mohīm yaśasvī jhālī.}
\end{quote}

\begin{quote}
The climbers' expedition was a success.
\end{quote}

Feminine.

\mr{मोहीम} \mr{आखणे}, to plan a campaign.

\begin{cousinweb}
\textbf{\mr{मोहीम}} (\emph{mohīm}) is a loan from Persian \emph{muhimm}, 'an important undertaking', usually traced to Arabic \emph{muhimm}, 'important'.
\end{cousinweb}

\begin{grammarlens}[title={What you've built}]
One more word for telling what happened.
\end{grammarlens}

\subsection*{Your turn}
\begin{itemize}
\raggedright
  \item \emph{Say it:} \emph{mohīm}
  \item \emph{Say it:} \emph{mohīm}, once more
  \item \emph{Say it:} \emph{giryārohkāñcī mohīm yaśasvī jhālī.} — The climbers' expedition was a success.
  \item \emph{Recall:} say the Marathi for a sacrifice, then the Marathi for to surrender, then say \emph{mohīm} again
\end{itemize}

\begin{tcolorbox}[breakable,colback=teal!4,colframe=teal!35!black,title={Before you move on}]
What does \mr{मोहीम} mean? (\textquotedblleft{}An expedition; a campaign\textquotedblright{}.) Say it once more.
\end{tcolorbox}
\section[\texorpdfstring{\mr{स्वारी} (svārī)}{स्वारी (svārī)}]{\mr{स्वारी} (svārī) — a military campaign; a ride}

\label{lesson:MR-C333-svari}



\begin{quote}
Before the new one: say the Marathi for to surrender, then the Marathi for an expedition.
\end{quote}

\subsection*{You'll want to know: \mr{स्वारी}}
\textbf{\mr{स्वारी}} — \emph{svārī} — \textquotedblleft{}a military campaign; a ride\textquotedblright{}.

\begin{quote}
\mr{राजाने} \mr{उत्तरेकडे} \mr{स्वारी} \mr{केली}.
\end{quote}

\begin{quote}
\emph{rājāne uttarekaḍe svārī kelī.}
\end{quote}

\begin{quote}
The king led a campaign to the north.
\end{quote}

Feminine.

\mr{स्वारी} \mr{करणे}, to march on a place.

\begin{cousinweb}
\textbf{\mr{स्वारी}} (\emph{svārī}) is a loan from Persian \emph{savārī}, 'riding, a mounted expedition', built on \emph{savār}, 'a rider'.
\end{cousinweb}

\begin{grammarlens}[title={What you've built}]
One more word for telling what happened.
\end{grammarlens}

\subsection*{Your turn}
\begin{itemize}
\raggedright
  \item \emph{Say it:} \emph{svārī}
  \item \emph{Say it:} \emph{svārī}, once more
  \item \emph{Say it:} \emph{rājāne uttarekaḍe svārī kelī.} — The king led a campaign to the north.
  \item \emph{Recall:} say the Marathi for to surrender, then the Marathi for an expedition, then say \emph{svārī} again
\end{itemize}

\begin{tcolorbox}[breakable,colback=teal!4,colframe=teal!35!black,title={Before you move on}]
What does \mr{स्वारी} mean? (\textquotedblleft{}A military campaign; a ride\textquotedblright{}.) Say it once more.
\end{tcolorbox}
\section[\texorpdfstring{\mr{माघार} \mr{घेणे} (māghār gheṇe)}{माघार घेणे (māghār gheṇe)}]{\mr{माघार} \mr{घेणे} (māghār gheṇe) — to retreat}

\label{lesson:MR-C333-maghar-ghene}



\begin{quote}
Before the new one: say the Marathi for an expedition, then the Marathi for a military campaign.
\end{quote}

\subsection*{You'll want to know: \mr{माघार} \mr{घेणे}}
\textbf{\mr{माघार} \mr{घेणे}} — \emph{māghār gheṇe} — \textquotedblleft{}to retreat\textquotedblright{}.

\begin{quote}
\mr{पावसामुळे} \mr{शत्रूने} \mr{माघार} \mr{घेतली}.
\end{quote}

\begin{quote}
\emph{pāvsāmuḷe śatrūne māghār ghetlī.}
\end{quote}

\begin{quote}
Because of the rain the enemy retreated.
\end{quote}

Past, with -\mr{ने} (he …): \mr{त्याने} \mr{माघार} \mr{घेतली}, \emph{tyāne māghār ghetlī}.

\mr{माघार} is feminine, so the past is \mr{माघार} \mr{घेतली}.

\begin{cousinweb}
\textbf{\mr{माघार} \mr{घेणे}} (\emph{māghār gheṇe}) joins \emph{māghār}, 'a falling back', to \emph{gheṇe}, 'to take'; the origin of \emph{māghār} is not settled.
\end{cousinweb}

\begin{grammarlens}[title={What you've built}]
One more word for telling what happened.
\end{grammarlens}

\subsection*{Your turn}
\begin{itemize}
\raggedright
  \item \emph{Say it:} \emph{māghār gheṇe}
  \item \emph{Say it:} \emph{māghār gheṇe}, once more
  \item \emph{Say it:} \emph{pāvsāmuḷe śatrūne māghār ghetlī.} — Because of the rain the enemy retreated.
  \item \emph{Recall:} say the Marathi for an expedition, then the Marathi for a military campaign, then say \emph{māghār gheṇe} again
\end{itemize}

\begin{tcolorbox}[breakable,colback=teal!4,colframe=teal!35!black,title={Before you move on}]
What does \mr{माघार} \mr{घेणे} mean? (\textquotedblleft{}To retreat\textquotedblright{}.) Say it once more.
\end{tcolorbox}
\section[\texorpdfstring{\mr{मुक्त} \mr{करणे} (mukta karṇe)}{मुक्त करणे (mukta karṇe)}]{\mr{मुक्त} \mr{करणे} (mukta karṇe) — to set free}

\label{lesson:MR-C333-mukta-karne}



\begin{quote}
Before the new one: say the Marathi for a military campaign, then the Marathi for to retreat.
\end{quote}

\subsection*{You'll want to know: \mr{मुक्त} \mr{करणे}}
\textbf{\mr{मुक्त} \mr{करणे}} — \emph{mukta karṇe} — \textquotedblleft{}to set free\textquotedblright{}.

\begin{quote}
\mr{राजाने} \mr{सगळ्या} \mr{कैद्यांना} \mr{मुक्त} \mr{केलं}.
\end{quote}

\begin{quote}
\emph{rājāne sagaḷyā kaidyānnā mukta kelaṁ.}
\end{quote}

\begin{quote}
The king set all the prisoners free.
\end{quote}

Past, with -\mr{ने} (he …): \mr{त्याने} \mr{मुक्त} \mr{केलं}, \emph{tyāne mukta kelaṁ}.

The one freed takes -\mr{ला}; \mr{मुक्त} \mr{होणे}, to become free.

\begin{cousinweb}
\textbf{\mr{मुक्त} \mr{करणे}} (\emph{mukta karṇe}) joins \emph{mukta}, a tatsama taken from Sanskrit \emph{mukta}, 'loosed, set free' (Monier-Williams), to \emph{karṇe}, 'to do'.
\end{cousinweb}

\begin{grammarlens}[title={What you've built}]
One more word for telling what happened.
\end{grammarlens}

\subsection*{Your turn}
\begin{itemize}
\raggedright
  \item \emph{Say it:} \emph{mukta karṇe}
  \item \emph{Say it:} \emph{mukta karṇe}, once more
  \item \emph{Say it:} \emph{rājāne sagaḷyā kaidyānnā mukta kelaṁ.} — The king set all the prisoners free.
  \item \emph{Recall:} say the Marathi for a military campaign, then the Marathi for to retreat, then say \emph{mukta karṇe} again
\end{itemize}

\begin{tcolorbox}[breakable,colback=teal!4,colframe=teal!35!black,title={Before you move on}]
What does \mr{मुक्त} \mr{करणे} mean? (\textquotedblleft{}To set free\textquotedblright{}.) Say it once more.
\end{tcolorbox}
\section[\texorpdfstring{\mr{हाकलणे} (hākalṇe)}{हाकलणे (hākalṇe)}]{\mr{हाकलणे} (hākalṇe) — to drive out; to chase away}

\label{lesson:MR-C333-hakalne}



\begin{quote}
Before the new one: say the Marathi for to retreat, then the Marathi for to set free.
\end{quote}

\subsection*{You'll want to know: \mr{हाकलणे}}
\textbf{\mr{हाकलणे}} — \emph{hākalṇe} — \textquotedblleft{}to drive out; to chase away\textquotedblright{}.

\begin{quote}
\mr{लोकांनी} \mr{शेतातून} \mr{माकडांना} \mr{हाकललं}.
\end{quote}

\begin{quote}
\emph{lokānnī śetātūn mākḍānnā hākallaṁ.}
\end{quote}

\begin{quote}
People drove the monkeys out of the field.
\end{quote}

Past, with -\mr{ने} (he …): \mr{त्याने} \mr{हाकललं}, \emph{tyāne hākallaṁ}.

The one driven out takes -\mr{ला}; the place left takes -\mr{तून}.

\begin{cousinweb}
\textbf{\mr{हाकलणे}} (\emph{hākalṇe}) is a tadbhava: CDIAL lists Marathi \emph{hākalṇe}, 'to drive on or away', under Sanskrit \emph{hakkayati}, 'calls to' (CDIAL 13939).
\end{cousinweb}

\begin{grammarlens}[title={What you've built}]
One more word for telling what happened.
\end{grammarlens}

\subsection*{Your turn}
\begin{itemize}
\raggedright
  \item \emph{Say it:} \emph{hākalṇe}
  \item \emph{Say it:} \emph{hākalṇe}, once more
  \item \emph{Say it:} \emph{lokānnī śetātūn mākḍānnā hākallaṁ.} — People drove the monkeys out of the field.
  \item \emph{Recall:} say the Marathi for to retreat, then the Marathi for to set free, then say \emph{hākalṇe} again
\end{itemize}

\begin{tcolorbox}[breakable,colback=teal!4,colframe=teal!35!black,title={Before you move on}]
What does \mr{हाकलणे} mean? (\textquotedblleft{}To drive out; to chase away\textquotedblright{}.) Say it once more.
\end{tcolorbox}
